% =============================================================================
% RHINO CW Calibration Experiment — Preparation Document
% Based on Jordan Norris / Phil Bull requirements and RFSoC 4x2 platform
% Prepared for: Mbatshi Jerry Junior Mbulawa
% Date: 2026-03-15
% =============================================================================
\documentclass[11pt,a4paper]{article}

\usepackage[a4paper,top=2.5cm,bottom=2.5cm,left=2.5cm,right=2.5cm]{geometry}
\usepackage[T1]{fontenc}
\usepackage[utf8]{inputenc}
\usepackage{xcolor}
\usepackage{listings}
\usepackage{booktabs}
\usepackage{longtable}
\usepackage{tabularx}
\usepackage{array}
\usepackage{enumitem}
\usepackage{titlesec}
\usepackage{fancyhdr}
\usepackage{tcolorbox}
\usepackage{hyperref}
\usepackage{amsmath}
\usepackage{amssymb}
\usepackage{parskip}

\tcbuselibrary{skins,breakable}

% Colours
\definecolor{xilinxblue}{HTML}{005590}
\definecolor{rhinored}{HTML}{B41E1E}
\definecolor{codebg}{HTML}{F0F4F8}
\definecolor{darkgrey}{HTML}{333333}
\definecolor{notebg}{HTML}{FFF8E1}
\definecolor{noteborder}{HTML}{E09B00}
\definecolor{warnbg}{HTML}{FFF0F0}
\definecolor{warnborder}{HTML}{B41E1E}
\definecolor{okbg}{HTML}{F0FFF4}
\definecolor{okborder}{HTML}{2D8A55}
\definecolor{infobg}{HTML}{EBF4FA}
\definecolor{codeborder}{HTML}{C0CDD8}

\hypersetup{colorlinks=true,linkcolor=xilinxblue,urlcolor=xilinxblue}

\lstset{
  basicstyle=\ttfamily\small,
  keywordstyle=\color{xilinxblue}\bfseries,
  commentstyle=\color{darkgrey}\itshape,
  stringstyle=\color{rhinored},
  backgroundcolor=\color{codebg},
  frame=single,framerule=0.4pt,rulecolor=\color{codeborder},
  breaklines=true,keepspaces=true,showstringspaces=false,
  xleftmargin=4pt,xrightmargin=4pt,aboveskip=4pt,belowskip=4pt,
  columns=flexible,tabsize=2,
}

\titleformat{\section}{\large\bfseries\color{xilinxblue}}{\thesection.}{0.6em}{}
  [\color{xilinxblue}\rule{\linewidth}{0.6pt}]
\titleformat{\subsection}{\normalsize\bfseries\color{darkgrey}}{\thesubsection}{0.5em}{}
\titleformat{\subsubsection}{\normalsize\bfseries\color{rhinored}}{\thesubsubsection}{0.5em}{}
\titlespacing{\section}{0pt}{12pt}{5pt}
\titlespacing{\subsection}{0pt}{8pt}{3pt}
\titlespacing{\subsubsection}{0pt}{6pt}{2pt}

\pagestyle{fancy}\fancyhf{}
\fancyhead[L]{\small\color{darkgrey}RHINO CW Calibration --- Experiment Preparation}
\fancyhead[R]{\small\color{darkgrey}2026-03-15}
\fancyfoot[C]{\small\color{darkgrey}\thepage}
\renewcommand{\headrulewidth}{0.3pt}

\newcolumntype{L}[1]{>{\raggedright\arraybackslash}p{#1}}

\newenvironment{infobox}{\begin{tcolorbox}[enhanced,colback=infobg,colframe=xilinxblue,
  left=5pt,right=5pt,top=3pt,bottom=3pt,before upper={\small}]}{\end{tcolorbox}}
\newenvironment{notebox}{\begin{tcolorbox}[enhanced,colback=notebg,colframe=noteborder,
  left=5pt,right=5pt,top=3pt,bottom=3pt,
  before upper={\small\textbf{\color{noteborder}NOTE:\enspace}}]}{\end{tcolorbox}}
\newenvironment{warnbox}{\begin{tcolorbox}[enhanced,colback=warnbg,colframe=warnborder,
  left=5pt,right=5pt,top=3pt,bottom=3pt,
  before upper={\small\textbf{\color{warnborder}WARNING:\enspace}}]}{\end{tcolorbox}}
\newenvironment{okbox}{\begin{tcolorbox}[enhanced,colback=okbg,colframe=okborder,
  left=5pt,right=5pt,top=3pt,bottom=3pt,
  before upper={\small\textbf{\color{okborder}CONFIRMED:\enspace}}]}{\end{tcolorbox}}
\newenvironment{todobox}{\begin{tcolorbox}[enhanced,colback=warnbg,colframe=rhinored,
  left=5pt,right=5pt,top=3pt,bottom=3pt,
  before upper={\small\textbf{\color{rhinored}TODO / OPEN QUESTION:\enspace}}]}{\end{tcolorbox}}

% =============================================================================
\begin{document}
% =============================================================================

\begin{center}
{\color{xilinxblue}\rule{\linewidth}{2pt}}\\[6pt]
{\LARGE\bfseries\color{xilinxblue}RHINO CW Calibration Experiment}\\[4pt]
{\large\bfseries Preparation Document: What to Expect, What to Build, What to Measure}\\[4pt]
{\normalsize Prepared for Mbatshi Jerry Junior Mbulawa \quad --- \quad 2026-03-15}\\[6pt]
{\color{xilinxblue}\rule{\linewidth}{2pt}}
\end{center}

\tableofcontents
\newpage

% =============================================================================
\section{Purpose of This Document}
% =============================================================================

This document consolidates everything known so far about the Continuous Wave (CW)
absolute radiometer calibration experiment being planned with Jordan Norris and
Dr.\ Phil Bull. It is intended as a preparation and thinking tool --- to be read
\textit{before} any code is written or any hardware is connected --- so that the
experiment is well understood before implementation begins.

The document covers:
\begin{itemize}[leftmargin=2em,itemsep=3pt]
  \item What the experiment is trying to achieve and why
  \item What Jordan has specified (and what remains open)
  \item The hardware and software constraints imposed by the RFSoC 4x2 platform
  \item A detailed analytical breakdown of what needs to be measured
  \item The two implementation paths (DMA buffer vs.\ DDS) and their trade-offs
  \item Open questions that must be resolved before implementation
  \item A concrete checklist of tasks
\end{itemize}

% =============================================================================
\section{Scientific Motivation}
% =============================================================================

\subsection{What Is Absolute Radiometer Calibration?}

A radiometer measures the total power in a band of frequencies. For the RHINO
experiment, the science band is 60--85\,MHz, targeting the redshifted 21\,cm
hydrogen emission line. However, any measurement of received power is contaminated
by:

\begin{itemize}[leftmargin=2em,itemsep=3pt]
  \item \textbf{Receiver noise} --- the thermal noise floor of the analogue front-end
        and ADC
  \item \textbf{Gain variations} --- the frequency-dependent transfer function of the
        signal chain (cables, amplifiers, filters, ADC)
  \item \textbf{ADC non-linearity} --- distortion introduced by the digitisation process
  \item \textbf{Spectral leakage} --- artefacts from the FFT windowing and finite
        sample length
\end{itemize}

\textit{Absolute} calibration means converting raw ADC power readings (in ADC counts
or dBFS) into physical units (e.g.\ Kelvin or Watts) with a known, characterised
uncertainty. Without this, you cannot distinguish a real signal from a systematic
artefact.

\subsection{Why Use a CW Tone?}

A continuous wave (CW) tone --- a pure sinusoid at a known frequency and amplitude
--- is the simplest possible calibration source. Its properties are:

\begin{itemize}[leftmargin=2em,itemsep=3pt]
  \item \textbf{Known power}: the injected power is set by the DAC amplitude, which is
        software-controllable
  \item \textbf{Known frequency}: the tone sits at exactly one FFT/PFB bin (if the
        system is correctly synchronised)
  \item \textbf{No spectral leakage}: because the ADC and DAC share the same clock on
        the RFSoC, there is no frequency offset between the injected tone and the FFT
        bin grid --- this eliminates one of the major sources of calibration error in
        traditional setups
\end{itemize}

By sweeping the CW tone across the science band and recording the ADC's measured
power at each frequency, you build up a \textit{gain curve} --- the system's frequency
response --- which can then be used to correct science observations.

\subsection{Connection to the Paper}

The RASTI paper (Jordan Norris et al., \textit{Continuous Wave Absolute Radiometer
Calibration}) will present this method as a novel calibration technique for low-cost
RFSoC-based radiometers. The experiment described in this document produces the
data that goes into the Methods and Results sections of that paper.

% =============================================================================
\section{Jordan's Requirements (Interpreted)}
% =============================================================================

\subsection{What Has Been Confirmed}

Based on the conversation with Jordan (2026-03-14 to 2026-03-15), the following
requirements are confirmed:

\begin{center}
\begin{tabular}{L{4cm} L{9.5cm}}
\toprule
\textbf{Requirement} & \textbf{Detail} \\
\midrule
Target frequency range & 50--200\,MHz \\
Signal type            & Sine wave (CW tone), zero DC offset \\
Controllable parameters & Amplitude and frequency (both must be quickly changeable) \\
Phase control          & Not required for now; total power measurement assumed \\
Clock synchronisation  & ADC and DAC must share the same clock (no spectral leakage) \\
Frequency switching    & Speed TBD --- both fast (microseconds) and slow (milliseconds)
                          approaches need to be tested to find what works \\
Current platform       & rfsoc\_sam overlay (custom bitstream not yet available) \\
\bottomrule
\end{tabular}
\end{center}

\subsection{Jordan's Final Steer: Flexibility First}

Following further discussion (2026-03-15), Jordan's overarching instruction is:

\begin{infobox}
\textbf{``There will need to be a lot of experimentation with this so I think a good
initial approach is to make things as flexible as possible.''}
\end{infobox}

This has a direct impact on the software architecture. Rather than optimising for
either fast or slow switching from the outset, the code must be written in a
\textit{backend-agnostic} way: the tone generation layer (NCO or DDS) must be
completely separable from the measurement layer (spectrum capture, averaging,
logging). Changing the tone generation method should require modifying only one
function, with zero impact on the rest of the pipeline.

\subsection{What Remains Open}

\begin{todobox}
\textbf{Open Question 2 --- Phase coherence:} Jordan said phase probably will not
matter for total power measurements, but may matter for some parts of calibration.
This is deferred for now. If phase coherence is later required, a DDS Compiler
implementation (see Section~\ref{sec:impl}) natively supports phase offset control
via AXI-Lite register writes.
\end{todobox}

\begin{todobox}
\textbf{Open Question 3 --- Absolute power reference:} To convert dBFS readings to
physical power (Watts or Kelvin), we need to know the injected power in absolute
terms. This requires either: (a) a calibrated attenuator or power meter in the
signal path, or (b) characterisation of the DAC output power as a function of
amplitude setting. This has not been discussed yet and must be resolved before the
paper can claim ``absolute'' calibration.
\end{todobox}

\begin{todobox}
\textbf{Open Question 4 --- Calibration metric:} What does a successful calibration
look like quantitatively? Possible answers include: the gain curve is flat to within
$X$\,dB, the noise figure matches a vendor spec, or the measured sky temperature
matches a known model. This needs to be defined with Jordan and Phil before the
experiment is run.
\end{todobox}

% =============================================================================
\section{Hardware Constraints: rfsoc\_sam in DDC Mode}
% =============================================================================

\subsection{The DDC Lock Problem}

As established in Session 1 of the RFSoC project (2026-03-04), the rfsoc\_sam
overlay is permanently configured in DDC (Digital Down Conversion) mode with the
ADC NCO fixed at $-1228.8$\,MHz. This cannot be changed at runtime via Python ---
the xrfdc C driver function \texttt{XRFdc\_MixerRangeCheck} blocks any attempt to
switch the mixer type.

This means that a signal at, say, 75\,MHz applied directly to the ADC\_A SMA
will be completely rejected by the DDC hardware filter. The signal must be
frequency-shifted before it enters the ADC.

\subsection{The Workaround: DAC NCO Mapping}

The solution is to use the RFSoC's own DAC to inject the CW tone at the
\textit{frequency-shifted} equivalent. The DDC offset formula is:

\begin{equation}
  f_{\text{DAC}} = 1228.8 + f_{\text{target}} \quad \text{(MHz)}
  \label{eq:ddc}
\end{equation}

where $f_{\text{target}}$ is the desired tone frequency in the ADC output spectrum.
For Jordan's range of 50--200\,MHz, this maps to DAC NCO values of:

\begin{equation}
  f_{\text{DAC}} \in [1278.8,\ 1428.8] \quad \text{MHz}
\end{equation}

This is well within the RFSoC DAC's RF bandwidth (up to 9.85\,GSPS on tile 228).

\subsection{Clock Synchronisation (Already Satisfied)}

Jordan's requirement for ADC/DAC clock synchronisation is automatically satisfied
on the RFSoC 4x2 platform. Both the ADC tile (Tile 224) and the DAC tile (Tile 228)
are clocked from the same on-chip PLL, which derives all sampling clocks from a
single reference. There is no external frequency offset between the two --- the tone
generated by the DAC will sit exactly on an FFT bin boundary in the ADC output
spectrum, provided the FFT length is chosen correctly (see Section~\ref{sec:fft}).

\subsection{FFT Bin Alignment Condition}
\label{sec:fft}

For a tone at $f_{\text{target}}$ to fall exactly on an FFT bin (zero spectral
leakage), the following condition must hold:

\begin{equation}
  f_{\text{target}} = k \cdot \frac{f_s}{N}
  \quad \text{for some integer } k
\end{equation}

where $f_s$ is the ADC output sample rate (200\,MSPS after 16$\times$ decimation in
rfsoc\_sam) and $N$ is the FFT length. With $f_s = 200\,\text{MSPS}$ and
$N = 4096$:

\begin{equation}
  \Delta f_{\text{bin}} = \frac{200}{4096} \approx 48.8 \quad \text{kHz}
\end{equation}

The tone frequency must therefore be a multiple of 48.8\,kHz for perfect bin
alignment. In practice, this means choosing $f_{\text{target}}$ values that are
multiples of this bin spacing when setting up the sweep. Non-aligned tones will
exhibit leakage even on a synchronised system.

\begin{notebox}
For a 1024-point FFT the bin spacing is 195.3\,kHz; for 8192 points it is 24.4\,kHz.
The choice of $N$ trades off frequency resolution against sensitivity. For calibration
purposes, finer resolution (larger $N$) is generally preferable.
\end{notebox}

% =============================================================================
\section{Software Architecture: Backend-Agnostic Design}
% =============================================================================

Given Jordan's flexibility requirement, the code is structured in three independent
layers, each of which can be modified or swapped without touching the others:

\begin{center}
\begin{tabular}{L{3cm} L{4cm} L{6.5cm}}
\toprule
\textbf{Layer} & \textbf{Responsibility} & \textbf{Swappable for} \\
\midrule
\textbf{Tone Generator} & Sets frequency, amplitude on the DAC &
  NCO (rfsoc\_sam, now) $\to$ DDS Compiler (custom bitstream, later) \\
\textbf{Spectrum Capture} & Reads ADC spectrum, averages $N$ frames &
  rfsoc\_sam FFT $\to$ DMA raw samples + software FFT \\
\textbf{Results Logger} & Saves all measurements to \texttt{.npy} / \texttt{.csv} &
  Unchanged regardless of backend \\
\bottomrule
\end{tabular}
\end{center}

This means that when the custom bitstream becomes available, only the
\texttt{set\_tone()} function needs to be rewritten. All measurement, plotting,
and logging code remains identical.

% =============================================================================
\section{Implementation Paths}
\label{sec:impl}
% =============================================================================

\subsection{Option A: DMA Buffer Approach (rfsoc\_sam, available now)}

This approach uses the existing rfsoc\_sam overlay without modification. The DAC
NCO is set to $f_{\text{DAC}} = 1228.8 + f_{\text{target}}$ for each tone, and
the transmitter amplitude is set in software.

\begin{center}
\begin{tabular}{L{3.5cm} L{9.5cm}}
\toprule
\textbf{Property} & \textbf{Detail} \\
\midrule
Switching speed    & Limited by Python overhead + DAC PLL re-lock time.
                     Estimated 10--100\,ms per step (needs measurement). \\
Amplitude control  & \texttt{ol.radio.transmitter.channel\_X.dac\_block.MixerSettings}
                     or amplitude scaling in the transmitter object. \\
Frequency control  & Set \texttt{MixerSettings['Freq']} on the DAC tile. \\
Phase control      & \texttt{MixerSettings['PhaseOffset']} --- arbitrary, not coherent. \\
Spectral purity    & High --- NCO-generated tone is spectrally pure (no DAC quantisation
                     noise from a buffer). \\
Implementation risk & Low --- no new Vivado work required. \\
\bottomrule
\end{tabular}
\end{center}

\textbf{Key limitation:} Every time you change the NCO frequency, the DAC PLL may
need to re-lock, introducing a dead time. This dead time must be characterised
(see Section~\ref{sec:tests}).

\subsubsection{Python Pseudocode for Option A}

\begin{lstlisting}[language=Python]
import numpy as np, time
from rfsoc_sam import Overlay   # rfsoc_sam overlay

ol = Overlay('rfsoc_sam.bit')
rx = ol.radio.receiver.channel_22   # ADC_A confirmed in Session 1
tx_tile = ol.radio.transmitter.channel_X   # replace X with correct channel

# Confirmed DDC offset from Session 1
DDC_OFFSET_MHz = 1228.8

def set_tone(freq_target_mhz, amplitude_norm):
    """Set DAC NCO to produce tone at freq_target_mhz in ADC spectrum."""
    f_dac = DDC_OFFSET_MHz + freq_target_mhz
    tx_tile.dac_block.MixerSettings['Freq'] = f_dac
    tx_tile.dac_block.UpdateEvent(1)   # apply settings
    tx_tile.amplitude = amplitude_norm
    tx_tile.transmit_enable = True
    time.sleep(0.05)   # allow PLL to re-lock --- tune this value

def measure_power(n_frames=20):
    """Capture N frames and return mean power spectrum in dBFS."""
    spectra = []
    for _ in range(n_frames):
        rx.spectrum_type = 'log'
        spectra.append(np.array(rx.spectrum))
    return np.mean(spectra, axis=0)

# Frequency sweep: 50-200 MHz in steps aligned to FFT bin grid
N_FFT = 4096
FS = 200e6   # 200 MSPS
bin_hz = FS / N_FFT
targets_mhz = np.arange(50, 201, bin_hz / 1e6)

results = {}
for f in targets_mhz:
    set_tone(f, amplitude_norm=0.5)
    spectrum = measure_power(n_frames=20)
    results[f] = spectrum
    print(f"f={f:.3f} MHz  peak={spectrum.max():.1f} dBFS")
\end{lstlisting}

\subsection{Option B: DDS Compiler (requires custom bitstream)}

This approach replaces the DMA DAC path in \texttt{system\_overlay} with a Xilinx
DDS Compiler IP block. Frequency, amplitude, and phase are all controlled via
AXI-Lite register writes, with switching latency in the microsecond range.

\begin{center}
\begin{tabular}{L{3.5cm} L{9.5cm}}
\toprule
\textbf{Property} & \textbf{Detail} \\
\midrule
Switching speed    & $\sim 1$\,$\mu$s (single AXI-Lite register write). \\
Amplitude control  & DDS output $\times$ AXI-Lite scale register. \\
Frequency control  & Phase increment register: $\Delta\phi = f_{\text{target}} \cdot 2^{32} / f_s$. \\
Phase control      & Phase offset register (coherent, relative to DDS reset). \\
Spectral purity    & Very high --- hardware DDS has controlled spurious levels. \\
Implementation risk & Medium --- requires Vivado work and new bitstream. \\
Licence dependency & Requires synthesis licence (currently blocked). \\
\bottomrule
\end{tabular}
\end{center}

\begin{warnbox}
Option B cannot be implemented until the Vivado floating licence tunnel to
\texttt{digdev4} (130.88.9.148) is resolved. Contact Ant
(\texttt{anthony.holloway@manchester.ac.uk}) for the jump host account on
\texttt{external.jb.man.ac.uk}, or use the University VPN.
\end{warnbox}

\subsection{Recommended Strategy}

Given Jordan's instruction to test both switching speeds, the recommended approach is:

\begin{enumerate}[leftmargin=2em,itemsep=4pt]
  \item \textbf{Start with Option A} (rfsoc\_sam). Measure the actual switching time
        experimentally (see Section~\ref{sec:tests}). If the switching time is fast
        enough for the calibration measurement, Option B may not be needed at all.
  \item \textbf{If Option A is too slow}, proceed to Option B once the bitstream
        licence is resolved. The DDS Compiler block can be added to
        \texttt{system\_overlay} without changing the ADC path.
\end{enumerate}

% =============================================================================
\section{Measurements to Make}
\label{sec:tests}
% =============================================================================

This section lists, in priority order, every measurement needed before the paper
can be written. Each measurement is annotated with whether it can be done now
(rfsoc\_sam) or requires the custom bitstream.

\subsection{Measurement 1: Switching Speed Characterisation (Option A)}
\label{sec:m1}

\textbf{What:} Measure the time from when the Python command to change the DAC NCO
frequency is issued to when the ADC spectrum shows a stable peak at the new
frequency.

\textbf{Why:} This determines whether Option A is fast enough for the calibration
workflow, and sets the minimum dwell time per frequency step.

\textbf{How:}
\begin{enumerate}[leftmargin=2em,itemsep=2pt]
  \item Set the DAC to $f_1 = 75$\,MHz in the ADC spectrum.
  \item Record timestamps using \texttt{time.perf\_counter()} before and after the
        NCO update command.
  \item Poll the ADC spectrum in a tight loop; record the time at which the peak
        at $f_2 = 100$\,MHz (the new frequency) first appears and stabilises
        (three consecutive frames with peak within 1\,dB of expected value).
  \item Repeat 20 times to get a distribution of switching times.
\end{enumerate}

\textbf{Platform:} rfsoc\_sam --- available now.

\textbf{Expected outcome:} Switching time of 10--100\,ms. If $< 10$\,ms, Option A
is likely sufficient. If $> 100$\,ms, Option B should be prioritised.

\subsection{Measurement 2: Gain Curve Sweep (50--200\,MHz)}
\label{sec:m2}

\textbf{What:} Sweep the CW tone from 50\,MHz to 200\,MHz in steps equal to one
FFT bin width ($\approx 48.8$\,kHz for $N = 4096$). At each step, record the
peak power in dBFS from the ADC spectrum.

\textbf{Why:} This is the primary calibration product --- the frequency-dependent
gain of the signal chain from DAC output to ADC FFT bin.

\textbf{How:}
\begin{enumerate}[leftmargin=2em,itemsep=2pt]
  \item Connect DAC\_A SMA to ADC\_A SMA via a known RF cable (record cable type
        and estimated insertion loss).
  \item For each target frequency $f$ (50 to 200\,MHz, bin-aligned):
        \begin{enumerate}[leftmargin=2em]
          \item Set DAC NCO to $1228.8 + f$\,MHz.
          \item Wait for PLL re-lock (dwell time from Measurement 1).
          \item Average 20 spectral frames.
          \item Record peak bin power and peak bin frequency.
        \end{enumerate}
  \item Store all results as a numpy array. Plot gain curve (dBFS vs.\ MHz).
\end{enumerate}

\textbf{Platform:} rfsoc\_sam --- available now.

\textbf{Expected outcome:} A smooth gain curve with $\pm$2--5\,dB variation across
the band. Large ripples ($> 5$\,dB) indicate cable reflections or connector issues.
A flat gain curve within the science band (60--85\,MHz) is the primary success
criterion.

\subsection{Measurement 3: Amplitude Linearity}
\label{sec:m3}

\textbf{What:} At a fixed frequency (e.g.\ 75\,MHz), sweep the DAC amplitude
from 0.1 to 1.0 (normalised) in 10 steps. Record the peak ADC power at each step.

\textbf{Why:} To verify that the ADC response is linear in power. If the relationship
between DAC amplitude and ADC peak power is linear in dB, the system is well-behaved
and the gain curve from Measurement 2 is amplitude-independent. Non-linearity
indicates ADC saturation or front-end compression.

\textbf{Expected outcome:} Peak power increases by 20\,dB for every factor-of-10
increase in amplitude (i.e.\ 6\,dB per factor-of-2 in amplitude). Any deviation
from this slope indicates non-linearity and must be characterised before calibration
results can be trusted.

\textbf{Platform:} rfsoc\_sam --- available now.

\subsection{Measurement 4: Noise Floor Stability}
\label{sec:m4}

\textbf{What:} With the DAC transmitter disabled (no tone injected), record the
ADC noise floor spectrum at 60, 70, 75, 80, and 85\,MHz over 10 minutes. Compute
the mean and standard deviation of the noise power in each bin over time.

\textbf{Why:} To characterise the radiometric stability of the system. For total
power calibration, you need to know how stable the noise floor is --- a drifting
noise floor means the calibration needs to be repeated frequently. This also sets
the minimum detectable signal level.

\textbf{Expected outcome:} Noise floor stable to within $\pm 0.5$\,dBFS over
10 minutes. Larger drifts indicate thermal effects or power supply instability.

\textbf{Platform:} rfsoc\_sam --- available now.

\subsection{Measurement 5: Switching Speed at Hardware Level (Option B)}
\label{sec:m5}

\textbf{What:} Repeat Measurement 1 using the DDS Compiler implementation
(AXI-Lite register write to change phase increment).

\textbf{Why:} To quantify the improvement in switching speed compared to Option A
and determine whether it is necessary for the calibration workflow.

\textbf{Platform:} Requires custom bitstream --- \textbf{blocked on licence}.

% =============================================================================
\section{What the Paper Will Need From These Measurements}
% =============================================================================

The RASTI paper will require the following from the experimental results:

\begin{center}
\begin{longtable}{L{3.5cm} L{4cm} L{5cm}}
\toprule
\textbf{Paper section} & \textbf{Required result} & \textbf{Source measurement} \\
\midrule
Methods & Hardware setup description & RFSoC 4x2 platform, rfsoc\_sam configuration \\
Methods & Signal generation method & Option A or B (to be decided after Measurement 1) \\
Methods & Calibration procedure & Measurements 2 and 3 \\
Results & Gain curve (50--200\,MHz) & Measurement 2 \\
Results & Amplitude linearity & Measurement 3 \\
Results & System noise floor & Measurement 4 \\
Results & Switching speed comparison & Measurements 1 and 5 \\
Conclusions & Calibration uncertainty & Derived from all measurements \\
\bottomrule
\end{longtable}
\end{center}

% =============================================================================
\section{Task Checklist}
% =============================================================================

\subsection{Immediate (This Week --- No Licence Needed)}

\begin{itemize}[leftmargin=2em,itemsep=4pt]
  \item[$\square$] Connect DAC\_A SMA to ADC\_A SMA on the board (loopback cable)
  \item[$\square$] Verify \texttt{channel\_22} is still the active ADC channel
        (re-run Session 1 loopback test if in doubt)
  \item[$\square$] Run Measurement 1: NCO switching speed characterisation
  \item[$\square$] Run Measurement 3: Amplitude linearity at 75\,MHz
  \item[$\square$] Run Measurement 4: Noise floor stability (10 min)
  \item[$\square$] Based on Measurement 1 result, decide with Jordan whether
        Option A is sufficient or Option B is needed
  \item[$\square$] Run Measurement 2: Full gain curve sweep (50--200\,MHz)
\end{itemize}

\subsection{Once Licence is Resolved}

\begin{itemize}[leftmargin=2em,itemsep=4pt]
  \item[$\square$] Generate bitstream for \texttt{system\_overlay}
  \item[$\square$] If Option B is needed: add DDS Compiler to \texttt{system\_overlay}
        and regenerate bitstream
  \item[$\square$] Run Measurement 5: DDS switching speed
  \item[$\square$] Re-run Measurements 2--4 with custom bitstream for comparison
\end{itemize}

\subsection{Before Writing the Paper}

\begin{itemize}[leftmargin=2em,itemsep=4pt]
  \item[$\square$] Resolve Open Question 3: define absolute power reference
  \item[$\square$] Resolve Open Question 4: define quantitative success criterion
        with Jordan and Phil
  \item[$\square$] Write Introduction (can be started now --- does not depend on
        experimental results)
\end{itemize}

% =============================================================================
\section{Key Numbers Summary}
% =============================================================================

\begin{center}
\begin{tabular}{L{6cm} L{7cm}}
\toprule
\textbf{Parameter} & \textbf{Value} \\
\midrule
Science band                & 60--85\,MHz \\
CW sweep range (Jordan)     & 50--200\,MHz \\
ADC sample rate (rfsoc\_sam) & 200\,MSPS (after 16$\times$ decimation) \\
DDC NCO offset              & 1228.8\,MHz \\
DAC NCO range for sweep     & 1278.8--1428.8\,MHz \\
FFT bin spacing ($N=4096$)  & $\approx 48.8$\,kHz \\
FFT bin spacing ($N=8192$)  & $\approx 24.4$\,kHz \\
Active ADC channel          & \texttt{channel\_22} (Tile 224, ADC 0) \\
Confirmed SNR (Session 1)   & 89.6\,dB at 100\,MHz offset \\
Noise floor (rfsoc\_sam)    & $-107.2$\,dBFS \\
\bottomrule
\end{tabular}
\end{center}

% =============================================================================
\section*{Acknowledgements}
% =============================================================================

This preparation document draws on work carried out with the support of
Dr.\ Phil Bull (University of Manchester / University of Western Cape),
Jordan Norris (University of Manchester), and Vladimir Stolyarov (Jodrell Bank
Observatory). The RFSoC 4x2 platform is provided by Real Digital. Vivado synthesis
tools are provided by AMD Xilinx.

\vspace{1.5cm}
\begin{center}
  \color{darkgrey}\rule{0.5\linewidth}{0.4pt}\\[4pt]
  {\small End of Preparation Document \quad --- \quad University of Manchester,
  RHINO RFSoC Project \quad --- \quad 2026-03-15}
\end{center}

\end{document}
